\documentclass[12pt]{article}
\usepackage{amsmath, amssymb}
\usepackage[margin=1in]{geometry}
\setlength{\parskip}{6pt}
\title{QUBO Formulation: Ligand Pose Hotspot Sampling}
\date{}
\begin{document}
\maketitle

\subsection*{Ligand Pose Hotspot Sampling}

\paragraph{Problem Statement.}
Let $V = \{0, 1, \dots, N-1\}$ denote the set of $N$ grid points inside a protein binding pocket. Each grid point $i \in V$ is associated with a three-dimensional coordinate $\mathbf{r}_i \in \mathbb{R}^3$ and an interaction energy score $E_i \in \mathbb{R}$. The Euclidean distance between any two distinct grid points $i$ and $j$ is denoted by $d_{ij} = \|\mathbf{r}_i - \mathbf{r}_j\|_2$. Given a minimum clash distance $d_{\text{clash}} \in \mathbb{R}_{>0}$, a preferred distance interval $[d_{\text{pref}}^{\min}, d_{\text{pref}}^{\max}] \subset \mathbb{R}_{\ge 0}$, a clash penalty strength $P_{\text{clash}} \in \mathbb{R}_{\ge 0}$, and a preferred-spacing reward strength $P_{\text{pref}} \in \mathbb{R}_{\ge 0}$, the task is to select a subset of grid points that minimizes the total interaction energy while penalizing spatial clashes and rewarding geometrically favorable separations. Formally, we define the clash pair set $C = \{(i,j) \in V \times V : i < j, d_{ij} < d_{\text{clash}}\}$ and the preferred pair set $P = \{(i,j) \in V \times V : i < j, d_{\text{pref}}^{\min} \le d_{ij} \le d_{\text{pref}}^{\max}\}$. The optimization problem seeks to minimize the combined energetic and geometric cost over all possible subsets.

\paragraph{Decision Variables.}
For each grid point $i \in V$, we introduce a binary decision variable $x_i \in \{0,1\}$. The variable $x_i$ indicates whether grid point $i$ is selected as a ligand-placement hotspot, where $x_i = 1$ denotes selection and $x_i = 0$ denotes non-selection. The decision vector is $\mathbf{x} = (x_0, x_1, \dots, x_{N-1})^\top \in \{0,1\}^N$.

\paragraph{QUBO Derivation.}
\textbf{Step 1.} The original objective function is a minimization problem that sums the interaction energies of selected points, adds quadratic penalties for selected pairs that are too close, and subtracts quadratic rewards for selected pairs within the preferred distance range. Since the QUBO framework natively minimizes quadratic forms, no sign flip is required. The objective is written as:
\begin{align}
    \min_{\mathbf{x} \in \{0,1\}^N} \left( \sum_{i \in V} E_i x_i + \sum_{(i,j) \in C} P_{\text{clash}} x_i x_j - \sum_{(i,j) \in P} P_{\text{pref}} x_i x_j \right).
\end{align}

\textbf{Step 2.} The problem specification contains no hard constraints; the subset size is free, and all spatial requirements are encoded as soft penalties/rewards in the objective. Thus, there are no constraint equations to enumerate.

\textbf{Step 3.} Because there are no hard constraints, no penalty terms are introduced. The Hamiltonian $H(\mathbf{x})$ is identical to the objective function derived in Step 1.

\textbf{Step 4.} Combining the objective and the (absent) penalty terms yields the QUBO Hamiltonian:
\begin{align}
    H(\mathbf{x}) = \sum_{i \in V} E_i x_i + \sum_{(i,j) \in C} P_{\text{clash}} x_i x_j - \sum_{(i,j) \in P} P_{\text{pref}} x_i x_j.
\end{align}

\textbf{Step 5.} Reading off the coefficients from the quadratic polynomial $H(\mathbf{x})$, we identify the symmetric QUBO matrix $Q$ and the constant offset $c$. The diagonal entries correspond to the linear coefficients, and the off-diagonal entries correspond to the quadratic coefficients. The general closed-form expressions are:
\begin{align}
    Q_{i,i} &= E_i, & \forall i \in V, \\
    Q_{i,j} &= \begin{cases}
        P_{\text{clash}} & \text{if } (i,j) \in C, \\
        -P_{\text{pref}} & \text{if } (i,j) \in P, \\
        0 & \text{otherwise},
    \end{cases} & \forall i,j \in V, i \neq j, \\
    c &= 0.
\end{align}
The QUBO is then expressed as $\min_{\mathbf{x} \in \{0,1\}^N} \left( \mathbf{x}^\top Q \mathbf{x} + c \right)$.

\paragraph{Penalty Weight Justification.}
Since the problem formulation contains no hard constraints, no constraint penalty weight $\lambda$ is required. The spatial requirements are handled entirely through the soft penalty and reward strengths $P_{\text{clash}}$ and $P_{\text{pref}}$ embedded directly in the objective. If hard constraints were introduced in a variant of this problem, the penalty weight $\lambda$ would be chosen such that $\lambda > \max_{\mathbf{x} \in \{0,1\}^N} |H_{\text{obj}}(\mathbf{x}) - H_{\text{obj}}(\mathbf{x}^*)|$, where $H_{\text{obj}}$ is the unconstrained objective and $\mathbf{x}^*$ is the optimal feasible solution. This ensures that any violation of a hard constraint increases the total energy by more than the maximum possible improvement achievable by any feasible configuration, thereby strictly penalizing infeasible states. For the current soft-constraint formulation, $P_{\text{clash}}$ and $P_{\text{pref}}$ are typically scaled relative to the magnitude of the energy scores $\{E_i\}_{i \in V}$ to maintain a balanced trade-off between energetic attraction and geometric plausibility.

\paragraph{Correctness Argument.}
Minimizing $H(\mathbf{x})$ over the Boolean hypercube $\{0,1\}^N$ directly recovers the optimal subset of grid points because the Hamiltonian is an exact algebraic representation of the problem's objective function. Since all requirements are encoded as soft terms in a single quadratic polynomial, every configuration $\mathbf{x} \in \{0,1\}^N$ is feasible, and its energy value precisely reflects the combined energetic and geometric cost. The global minimum of $H(\mathbf{x})$ therefore corresponds exactly to the subset that minimizes total interaction energy while optimally balancing clash avoidance and preferred spacing, without any feasible/infeasible distinction or approximation.

\paragraph{Illustrative Example.}
Consider a concrete instance with $N=3$ grid points, so $V = \{0, 1, 2\}$. For demonstration, assume all distinct pairs fall into both the clash and preferred sets (i.e., $C = P = \{(0,1), (0,2), (1,2)\}$). We instantiate the general parameters with the following values: $E_0 = -5$, $E_1 = -3$, $E_2 = -4$, $P_{\text{clash}} = 10$, and $P_{\text{pref}} = 2$. Substituting these into the general objective yields the concrete energy function:
\begin{align}
    H(\mathbf{x}) &= -5x_0 - 3x_1 - 4x_2 + 10(x_0x_1 + x_0x_2 + x_1x_2) - 2(x_0x_1 + x_0x_2 + x_1x_2) \\
    &= -5x_0 - 3x_1 - 4x_2 + 8(x_0x_1 + x_0x_2 + x_1x_2).
\end{align}
The corresponding QUBO matrix entries and offset are:
\begin{align}
    Q_{0,0} &= -5, \quad Q_{1,1} = -3, \quad Q_{2,2} = -4, \\
    Q_{0,1} = Q_{1,0} &= 8, \quad Q_{0,2} = Q_{2,0} = 8, \quad Q_{1,2} = Q_{2,1} = 8, \\
    c &= 0.
\end{align}
Enumerating all $2^3 = 8$ binary configurations reveals the following energies: $H(1,0,0) = -5$, $H(0,1,0) = -3$, $H(0,0,1) = -4$, $H(1,1,0) = 0$, $H(1,0,1) = -1$, $H(0,1,1) = 1$, $H(1,1,1) = 12$, and $H(0,0,0) = 0$. The global minimum occurs at $\mathbf{x}^* = (1, 0, 0)^\top$ with an optimal energy of $-5$, indicating that only grid point $0$ should be selected as a hotspot for this instance.

\end{document}